\clearpage
\fancyhead[L]{\small GAUSS TO VON LINDENAU.}
\hypertarget{native-301-lower}{}
\begin{center}GAUSS TO VON LINDENAU.
\end{center}

\begin{center}
\begin{tabular}{c}
\hline
Monatliche Correspondenz. Vol. XVII. Pp. 182--188. February 1808.\\
\hline
\end{tabular}
\end{center}

\begin{flushright}
Göttingen, 25 January 1808.
\end{flushright}

Please excuse me for being unable to send you a series of good observations of the comet from here as I had hoped. The weather has been unprecedentedly bad. In all the time I have been here (since 21 November), there have been only 3 evenings on which I managed to wrest a few scant observations from the unfavorable sky. I give them here with the differences from my elements. [These are the observations of 17 and 29 December and 4 January printed below on page 298 in the reprint of the article of 25 February 1808 from the Göttingische gelehrte Anzeigen.]

\par\indent Less reliance can be placed on the observations of 29 December because the instrument used to make them is poorly mounted. The first two observations of 4 January are only isolated determinations made between clouds, but I consider them good in themselves. The last is by my esteemed colleague, Prof. \textsc{Harding}, and is the mean of several comparisons that agree well. In any case, the recent observations of the comet cannot be very accurate. It is so faint and its nucleus so hard to make out that its entries and exits can no longer be observed so precisely. Incidentally, the inclination of the orbit was omitted from my elements in the December issue of the M.\,C. [p. 292 of this volume]. It is $63^\circ\;12'\;35''6$. I see the omission in the copy belonging to Court Councillor \textsc{Blumenbach} that has just been shown to me.

\par\indent I am much obliged to you for kindly sending me the Milan observations of \VestaAltarThreeOhOne{}. My excellent friend Mr. \textsc{Bessel} was kind enough to reduce the latest ones with the greatest care, taking proper account of refraction. Here are the results of those reductions:


\begin{center}\fontsize{9}{11}\selectfont\setlength{\tabcolsep}{2pt}\renewcommand{\arraystretch}{1.15}\begin{tabular}{lrrrrr}
& & & & \multicolumn{2}{c}{Differences}\\
1807 & Mean time & R.A. & Declination & R.A. & Decl.\\\hline
Sept. 5 & $7^{\text{u}}\,21^{\text{m}}\,13^{\text{s}}0$ & $211^\circ\;29'\;4''5$ & $8^\circ\;2'\;49''9$ south & $-46''6$ & $-53''0$\\
8 & $7\;14\;47.5$ & $212\;49\;55.4$ & $8\;38\;18.1$ & $-54.0$ & $-59.3$\\
8 & $7\;29\;44.8$ & $212\;50\;11.5$ & $8\;38\;16.0$ & $-53.2$ & $-49.9$\\
10 & $7\;13\;14.9$ & $213\;44\;31.7$ & $9\;1\;45.8$ & $-65.5$ & $-61.4$\\
21 & $7\;4\;10.8$ & $218\;51\;35.2$ & $11\;7\;42.5$ & $-82.4$ & $-64.5$\\
26 & $6\;53\;51.8$ & $221\;15\;21.1$ & $12\;2\;59.3$ & $-99.5$ & $-66.9$\\
\end{tabular}\end{center}

\par\indent Prof. \textsc{Harding} kindly compared these observations with my III. elements. The comparison gave the following differences. [In this reprint they have been added to the observations above.]

\par\indent I\SourceLeaf{302} see no need to improve the III. elements again at this point on the basis of these observations. The corrections would be very small, and the ephemeris in the September issue of the M.\,C. will suffice to find the planet.

\par\indent Here are the new elements of \emph{Pallas} from which the ephemeris in the December issue of the M.\,C. was calculated:

\begin{center}
X$^{\text{th}}$ Elements
\end{center}


\begin{center}\fontsize{10}{12}\selectfont\setlength{\tabcolsep}{2pt}\renewcommand{\arraystretch}{1.15}\begin{tabular}{cr}
Epoch & \shortstack{Göttingen\\meridian}\\\hline
1802 & $143^\circ\;33'\;58''7$\\
1803 & $221\;39\;26.9$\\
1804 & $299\;57\;45.4$\\
1805 & $18\;3\;13.6$\\
1806 & $96\;8\;41.8$\\
1807 & $174\;14\;10.0$\\
1808 & $252\;32\;28.5$\\
\end{tabular}\end{center}

\begin{center}\fontsize{10}{12}\selectfont\setlength{\tabcolsep}{2pt}\renewcommand{\arraystretch}{1.15}\begin{tabular}{lr}
Daily tropical motion .\;.\;.\;.\;.\; & $770^{.}2143$\\
Logarithm of the semimajor axis .\;.\; & $0.4423149$\\
Eccentricity .\;.\;.\;.\;.\;.\;.\;.\;.\;.\;.\; & $0.2450198$\\
Perihelion 1803 .\;.\;.\;.\;.\;.\;.\; & $121^\circ\;3'\;11''4$\\
$\AscendingNode$, 1803 .\;.\;.\;.\;.\;.\;.\;.\;.\;.\; & $172\;28\;56.9$\\
Inclination of the orbit .\;.\;.\;.\;.\;.\;.\; & $34\;37\;41.0$\\
\end{tabular}\end{center}

\par\indent These elements are based solely on the four oppositions of 1803, 1804, 1805 and 1807. The observations of 1802 were left entirely out of account because they are unsuitable for determining the opposition precisely. All the longitudes of those oppositions are represented accurately,, and although the latitudes cannot be made to agree exactly, the differences fall short of $\frac12$ minute. Pallas will in any event be harder to observe in 1808 than ever, since it is now at aphelion.

\Needspace{19\baselineskip}
\par\indent Some time ago I also improved the elements of \emph{Juno}. You will find enclosed the results and the ephemeris for 1808 calculated from them. I will only add that Prof. \textsc{Harding} kindly compared the new VII$^{\text{th}}$ elements carefully with Mr. \textsc{Bessel}’s observations, leaving parallax out of account.


\begin{center}\fontsize{9}{11}\selectfont\setlength{\tabcolsep}{2pt}\renewcommand{\arraystretch}{1.15}\begin{tabular}{lrrrrr}
1807 & \shortstack{Mean time\\Lilienthal} & \shortstack{Calculated\\R.A.} & \shortstack{Diff.\\R--B} & \shortstack{Calc. south\\declination} & \shortstack{Diff.\\R--B}\\\hline
May 5 & $12^{\text{u}}\,46^{\text{m}}\,4^{\text{s}}$ & $240^\circ\;5'\;11''4$ & $+26''6$ & $3^\circ\;44'\;27''1$ & $+7''8$\\
7 & $12\;25\;33$ & $239\;42\;31.0$ & $+3.6$ & $3\;34\;18.0$ & $+17.6$\\
8 & $11\;30\;29$ & $239\;31\;22.3$ & $+5.6$ & $3\;29\;29.4$ & $+0.1$\\
19 & $12\;8\;26$ & $237\;17\;26.9$ & $-5.6$ & $2\;40\;34.5$ & $-3.8$\\
23 & $11\;51\;8$ & $236\;28\;5.2$ & $-9.5$ & $2\;26\;12.1$ & $+5.9$\\
24 & $11\;14\;57$ & $236\;16\;7.4$ & $-5.2$ & $2\;22\;59.9$ & $-0.2$\\
25 & $11\;54\;26$ & $236\;3\;31.0$ & $-2.6$ & $2\;19\;44.7$ & $-2.3$\\
26 & $11\;50\;58$ & $235\;51\;11.4$ & $-5.7$ & $2\;16\;41.9$ & $-16.0$\\
\end{tabular}\end{center}

\par\indent I have a misprint to correct in the ephemeris for $\PallasSymbol$. The last column gives the logarithms of the distances of Pallas from Earth. The reference to the distances themselves is a misprint.

\par\indent My work on determining the orbits of celestial bodies is now in press and will, I hope, be able to appear soon after Easter.

\par\indent I am also sending you some small tables of aberration and nutation. I calculated them four years ago and shared them with several friends, who found them convenient and have grown used to working with them. Of all the general tables, these seem to me to make calculations quickest and shortest. The need for tables of sines and logarithms hardly matters, since every astronomer already keeps them on their desk. The copy of th\SourceLeaf{303}ese tables I am sending you here was calculated afresh from Mr. \textsc{von Zach}’s latest numerical data. [Pages 123..128 of this volume.]

\par\indent [The geocentric path of Juno is given for midnight in Göttingen at four-day intervals from 16 April to 28 December 1808. Right ascension and declination are given to the nearest minute, and the logarithm of the distance from Earth to four decimal places.]

